\documentclass[../../main]{subfiles}

\begin{document}

\section{Composição de Funções}
\label{sec:matematica:funcoes:composicao}

\begin{defn}[Composição de Funções]
	\label{def:matematica:funcoes:composicao:composicao}

	Sejam

		[
			f:A\to B
		]

	e

		[
			g:B\to C.
		]

	A composição de (g) com (f) é a função

	[
	g\circ f:A\to C
	]

	definida por

		[
			(g\circ f)(x)
			=
			g(f(x))
		]

	para todo (x\in A).

\end{defn}

\begin{obs}

	Na expressão

	[
	g\circ f,
	]

	a função (f) é aplicada primeiro e a função (g) é aplicada em
	seguida.

\end{obs}

\begin{prop}[Boa Definição]
	\label{prop:matematica:funcoes:composicao:boa-definicao}

	Se

		[
			f:A\to B
		]

	e

		[
			g:B\to C,
		]

	então

	[
	g\circ f
	]

	é uma função de (A) em (C).

\end{prop}

\begin{proof}

	Seja (x\in A).

	Como

		[
			f:A\to B,
		]

	temos

		[
			f(x)\in B.
		]

	Como

		[
			g:B\to C,
		]

	existe um único elemento

		[
			g(f(x))
			\in
			C.
		]

	Portanto

		[
			g\circ f:A\to C
		]

	satisfaz as propriedades de existência e unicidade da imagem.

\end{proof}

\begin{prop}[Associatividade]
	\label{prop:matematica:funcoes:composicao:associatividade}

	Sejam

		[
			f:A\to B,
		]

		[
			g:B\to C
		]

	e

		[
			h:C\to D.
		]

	Então

	[
	h\circ(g\circ f)
	=
	(h\circ g)\circ f.
	]

\end{prop}

\begin{proof}

	Para todo (x\in A),

	[
			(h\circ(g\circ f))(x)
			=
			h((g\circ f)(x))
		]

		[
			=
			h(g(f(x))).
		]

	Por outro lado,

	[
			((h\circ g)\circ f)(x)
			=
			(h\circ g)(f(x))
		]

		[
			=
			h(g(f(x))).
		]

	Portanto as duas funções coincidem em todos os elementos do domínio.

	Pelo princípio da extensionalidade,

	[
			h\circ(g\circ f)
			=
			(h\circ g)\circ f.
		]

\end{proof}

\subsection{Funções Identidade}

\begin{defn}[Função Identidade]
	\label{def:matematica:funcoes:composicao:identidade}

	Seja (A) um conjunto.

	A função identidade de (A) é a função

	[
	\operatorname{id}_A:A\to A
	]

	definida por

		[
			\operatorname{id}_A(x)=x
		]

	para todo (x\in A).

\end{defn}

\begin{prop}[Neutro à Esquerda]
	\label{prop:matematica:funcoes:composicao:neutro-esquerda}

	Seja

		[
			f:A\to B.
		]

	Então

	[
	\operatorname{id}_B
	\circ
	f
	=
	f.
	]

\end{prop}

\begin{proof}

	Para todo (x\in A),

	[
			(\operatorname{id}_B\circ f)(x)
			=
			\operatorname{id}_B(f(x))
			=
			f(x).
		]

	Pelo princípio da extensionalidade,

	[
			\operatorname{id}_B\circ f=f.
		]

\end{proof}

\begin{prop}[Neutro à Direita]
	\label{prop:matematica:funcoes:composicao:neutro-direita}

	Seja

		[
			f:A\to B.
		]

	Então

	[
	f
	\circ
	\operatorname{id}_A
	=
	f.
	]

\end{prop}

\begin{proof}

	Para todo (x\in A),

	[
			(f\circ\operatorname{id}_A)(x)
			=
			f(\operatorname{id}_A(x))
			=
			f(x).
		]

	Pelo princípio da extensionalidade,

	[
			f\circ\operatorname{id}_A=f.
		]

\end{proof}

\subsection{Potências de Endofunções}

\begin{defn}[Potência de uma Endofunção]
	\label{def:matematica:funcoes:composicao:potencia}

	Seja

		[
			f:A\to A.
		]

	Define-se recursivamente

		[
			f^0
			=
			\operatorname{id}_A,
		]

	e

		[
			f^{n+1}
			=
			f\circ f^n.
		]

\end{defn}

\begin{ex}

	[
		f^1=f,
	]

	[
		f^2=f\circ f,
	]

	[
		f^3=f\circ f\circ f.
	]

\end{ex}

\begin{prop}[Lei das Potências]
	\label{prop:matematica:funcoes:composicao:lei-potencias}

	Para quaisquer

		[
			m,n\in\mathbb N,
		]

	vale

		[
			f^m\circ f^n
			=
			f^{m+n}.
		]

\end{prop}

\begin{proof}

	Procedemos por indução em (m).

	Para (m=0),

	[
			f^0\circ f^n
			=
			\operatorname{id}_A\circ f^n
			=
			f^n.
		]

	Portanto

		[
			f^0\circ f^n
			=
			f^{0+n}.
		]

	Suponha

		[
			f^m\circ f^n
			=
			f^{m+n}.
		]

	Então

	[
	f^{m+1}\circ f^n
	=
	(f\circ f^m)\circ f^n.
	]

	Pela associatividade,

	[
			=
			f\circ(f^m\circ f^n).
		]

	Pela hipótese de indução,

	[
			=
			f\circ f^{m+n}.
		]

	Logo

		[
			=
			f^{m+n+1}
			=
			f^{(m+1)+n}.
		]

\end{proof}

\subsection{Cancelamento}

\begin{prop}[Cancelamento à Esquerda]
	\label{prop:matematica:funcoes:composicao:cancelamento-esquerda}

	Seja

		[
			f:A\to B
		]

	uma função injetiva.

	Se

		[
			f\circ g
			=
			f\circ h,
		]

	então

	[
	g=h.
	]

\end{prop}

\begin{proof}

	Para todo elemento (x),

	[
			f(g(x))
			=
			f(h(x)).
		]

	Como (f) é injetiva,

	[
			g(x)=h(x).
		]

	Pelo princípio da extensionalidade,

	[
			g=h.
		]

\end{proof}

\begin{prop}[Cancelamento à Direita]
	\label{prop:matematica:funcoes:composicao:cancelamento-direita}

	Seja

		[
			f:A\to B
		]

	uma função sobrejetiva.

	Se

		[
			g\circ f
			=
			h\circ f,
		]

	então

	[
	g=h.
	]

\end{prop}

\begin{proof}

	Seja (y\in B).

	Como (f) é sobrejetiva, existe
	(x\in A) tal que

		[
			f(x)=y.
		]

	Logo

		[
			g(y)
			=
			g(f(x))
			=
			h(f(x))
			=
			h(y).
		]

	Pelo princípio da extensionalidade,

	[
			g=h.
		]

\end{proof}

\end{document}
